\documentclass[conference]{IEEEtran}

\usepackage{cite}
\usepackage{amsmath,amssymb,amsfonts}
\usepackage{algorithmic}
\usepackage{graphicx}
\usepackage{textcomp}
\usepackage{xcolor}
\usepackage{hyperref}
\usepackage{listings}
\usepackage{booktabs}
\usepackage{multirow}

\lstset{
  basicstyle=\ttfamily\small,
  breaklines=true,
  frame=single,
  numbers=left,
  numberstyle=\tiny,
}

\begin{document}

\title{AgentSentry: A Risk Quantification Framework for\\
Non-Human Identities and Autonomous AI Agents\\
in Cloud Environments}

\author{
  \IEEEauthorblockN{Abhiram Lanka}
  \IEEEauthorblockA{
    \textit{Independent Security Researcher}\\
    hello@agentsentry.org\\
    \href{https://github.com/Abhiram-ops/agent-sentry}{github.com/Abhiram-ops/agent-sentry}
  }
}

\maketitle

\begin{abstract}
Modern cloud environments contain roughly 45 non-human identities (NHIs) for
every human identity, yet existing identity security frameworks offer no
rigorous methodology for quantifying the risk these credentials actually pose.
The rise of autonomous AI agents makes this gap considerably worse. These
agents, which execute real-world actions without waiting for human approval,
introduce a new category of identity whose danger scales with autonomy just
as much as with privilege.

We present \textbf{AgentSentry}, an open-source auditing framework that
discovers, inventories, and risk-scores NHIs and AI agents across cloud and
local environments. The core of our contribution is the \textbf{AI-Amplification
Factor} ($A$), a novel scoring variable that captures how an autonomous AI agent
running under a compromised identity multiplies the damage that compromise can
cause. We formalize a four-component risk model $\text{Risk} = P \times R
\times E \times A$, build it into a working open-source tool, and test it
against live AWS environments and AI agent codebases. The results show the
$A$ factor pushes composite risk scores up to $60\times$ higher than equivalent
non-agentic identities, surfacing a class of critical vulnerabilities that
existing tools simply cannot see.
\end{abstract}

\begin{IEEEkeywords}
non-human identity, AI agents, cloud security, identity and access management,
attack graph, risk quantification, LangChain, autonomous agents
\end{IEEEkeywords}

\section{Introduction}

Cloud-native architectures have quietly changed the composition of enterprise
identity. Traditional IAM frameworks assumed that most authentication events
came from humans. That assumption no longer holds. Today's enterprise runs on
\emph{non-human identities} (NHIs): IAM roles, service accounts, API keys,
OAuth tokens, and CI/CD credentials that handle machine-to-machine
communication continuously in the background.

Industry measurements consistently put the ratio of machine identities to human
ones somewhere between 40:1 and 50:1 in large organizations
\cite{crowdstrike2024}. In spite of this, NHIs receive a fraction of the
security attention and tooling investment that human accounts do. Frameworks
such as NIST SP 800-53, CIS Controls, and various Cloud Security Posture
Management (CSPM) products treat NHI governance as a secondary concern, folded
into broader identity policy rather than addressed on its own terms.

Autonomous AI agents make this situation meaningfully worse. Platforms like
LangChain~\cite{langchain2023}, CrewAI~\cite{crewai2024}, and
AutoGen~\cite{autogen2023} let developers deploy LLM-powered agents that pick
and invoke \emph{tools} on their own: functions that write to databases, send
emails, execute code, or move money. When such an agent runs under an identity
that has been compromised or over-permissioned, the blast radius is not just
the set of resources that identity can reach. It extends to every irreversible
action the agent can take at machine speed, with no human in the loop to notice
something has gone wrong.

No current security framework puts a number on this amplification. MITRE
ATT\&CK~\cite{mitre2024} documents the techniques attackers use against AI
systems, but it does not produce a quantitative risk score for agentic identity
exposure. The OWASP Top 10 for LLM Applications~\cite{owasp_llm2023} tackles
prompt injection and insecure output handling, but stops short of modeling
identity-level blast radius. This paper addresses that gap directly.

Our contributions are as follows:

\begin{itemize}
  \item We define the \textbf{AI-Amplification Factor} ($A$), the first formal
        variable to quantify how autonomous AI agent execution compounds the risk
        of a compromised NHI.
  \item We present a four-component risk scoring model
        $\text{Risk} = P \times R \times E \times A$ that extends privilege-based
        scoring with reachability, credential lifecycle exposure, and agentic
        amplification.
  \item We implement the model in \textbf{AgentSentry}, an open-source Python
        tool that analyzes cloud IAM configurations and AI agent codebases,
        constructs an NHI attack graph, and produces prioritized findings mapped
        to MITRE ATT\&CK and enriched with CISA Known Exploited Vulnerability
        (KEV) data.
  \item We validate the framework against a live AWS environment and show that
        agentic NHIs score up to $60\times$ higher than non-agentic identities
        at the same privilege level.
\end{itemize}

\section{Background and Related Work}

\subsection{Non-Human Identity Governance}

The term \emph{non-human identity} refers to any credential or identity used by
a software process rather than a person. AWS IAM roles~\cite{aws_iam2024},
Azure managed identities~\cite{azure_identity2024}, and Kubernetes service
accounts~\cite{k8s_sa2024} are the most common examples. Secretless
authentication~\cite{conjur2022} and workload identity
federation~\cite{workload_identity2023} represent better-practice alternatives,
though adoption varies widely across organizations.

Prior work on NHI risk has focused almost entirely on privilege analysis. The
principle of least privilege~\cite{saltzer1975} provides the theoretical
foundation; tools such as AWS IAM Access Analyzer~\cite{access_analyzer} and
commercial CIEM platforms operationalize it. These approaches treat NHI risk as
a pure function of what permissions an identity holds. They do not account for
how long credentials have sat unrotated, whether they are accessible from the
internet, or what an AI agent running under them might do autonomously.

\subsection{AI Agent Security}

Serious attention to LLM agent security began emerging around 2023. Prompt
injection~\cite{perez2022} is the most studied attack vector: adversarial
content in an agent's input can override its intended behavior. Indirect prompt
injection~\cite{greshake2023}, where malicious instructions arrive through
documents retrieved in RAG pipelines rather than directly, is particularly
relevant to NHI security because it lets attackers influence agent tool calls
without ever touching the agent's primary input.

Yi et al.~\cite{yi2023} examine tool-use safety in LLM agents, with an
emphasis on output correctness. To our knowledge, no prior work formally models
how autonomous agent execution multiplies the blast radius of a compromised
NHI. That is the gap this work fills.

\subsection{Attack Graph Models}

Attack graphs~\cite{sheyner2002} have been used in network security for
decades as a formal way to reason about adversary reachability. More recent
research has extended attack graph techniques to cloud
environments~\cite{orca2023, wiz2023}. AgentSentry applies this approach to
the NHI domain, building a directed graph where nodes are identities and
resources and edges are access relationships, then using standard graph
reachability to compute blast radius automatically.

\section{Threat Model}

We model an adversary trying to maximize impact from a single successful
initial access event against a cloud environment. The specific assumptions are:

\begin{itemize}
  \item \textbf{Initial access vector}: Compromise of one NHI through credential
        theft, secret exposure, or IAM misconfiguration.
  \item \textbf{Objective}: Reach as many sensitive resources (crown jewels) as
        possible through as few lateral movement steps as possible.
  \item \textbf{Capability}: The adversary can use any permission the compromised
        identity holds, invoke any tool available to an AI agent running under
        that identity, and exploit prompt injection if a RAG pipeline is reachable.
  \item \textbf{Constraint}: The adversary works within the authorization boundary
        of the compromised identity. We do not assume zero-day exploitation of
        cloud provider infrastructure.
\end{itemize}

This model maps to MITRE ATT\&CK Cloud Matrix techniques including T1078.004
(Valid Accounts: Cloud Accounts), T1528 (Steal Application Access Token),
T1199 (Trusted Relationship), and T1651 (Cloud Administration Command).

\section{The AgentSentry Scoring Framework}

\subsection{Overview}

The \textbf{NHI Risk Score} is defined as a product of four components:

\begin{equation}
  \text{Risk}(i) = P(i) \times R(i) \times E(i) \times A(i)
  \label{eq:risk}
\end{equation}

where $i$ is a specific NHI, $P$ is the Privilege Score, $R$ is the
Reachability Score, $E$ is the Exposure Score, and $A$ is the AI-Amplification
Factor. We use multiplication rather than addition because these risk dimensions
compound. An identity that is highly privileged, externally reachable, and
never rotated, operated by an autonomous AI agent, presents risk that is the
product of all four conditions, not their sum.

Risk levels map to score ranges:
$\text{CRITICAL} \geq 100$,
$\text{HIGH} \in [50, 100)$,
$\text{MEDIUM} \in [20, 50)$,
$\text{LOW} < 20$.

\subsection{Privilege Score ($P$)}

The Privilege Score places an identity's access permissions on a 1-to-10 scale:

\begin{equation}
  P(i) = \min\!\left(10,\; \max_{a \in \text{Actions}(i)} w(a) \times
         \mathbb{1}[\text{cross-account}(i)] \cdot 1.5\right)
  \label{eq:privilege}
\end{equation}

Here, $w(a)$ is a weight reflecting the destructive potential of action $a$
(Table~\ref{tab:permission_weights}), and the cross-account multiplier accounts
for the added exposure of identities with external trust relationships.
Administrator-level policies set $P = 10$ unconditionally.

\begin{table}[h]
\caption{Permission Weight Table (selected entries)}
\label{tab:permission_weights}
\centering
\begin{tabular}{lc}
\toprule
\textbf{Permission} & \textbf{Weight} \\
\midrule
\texttt{iam:*} (wildcard) & 9.0 \\
\texttt{iam:AttachRolePolicy} & 8.5 \\
\texttt{iam:CreateUser} & 8.0 \\
\texttt{iam:PassRole} & 7.5 \\
\texttt{sts:AssumeRole} & 7.0 \\
\texttt{secretsmanager:PutSecretValue} & 6.0 \\
\texttt{lambda:UpdateFunctionCode} & 6.0 \\
\texttt{s3:PutBucketPolicy} & 6.0 \\
\texttt{s3:PutObject} & 2.5 \\
\texttt{s3:GetObject} (read-only) & 1.0 \\
\bottomrule
\end{tabular}
\end{table}

\subsection{Reachability Score ($R$)}

The Reachability Score captures how accessible an identity is to an external
attacker:

\begin{equation}
  R(i) = \begin{cases}
    3.0 & \text{if internet-facing} \\
    2.5 & \text{if AI agent type} \\
    2.0 & \text{if API key or GitHub secret} \\
    1.0 & \text{otherwise (internal only)}
  \end{cases}
  \label{eq:reachability}
\end{equation}

\subsection{Exposure Score ($E$)}

The Exposure Score penalizes poor credential lifecycle practices:

\begin{equation}
  E(i) = \min\!\left(5,\; f_{\text{rot}}(i) \times f_{\text{use}}(i)
         \times f_{\text{gov}}(i)\right)
  \label{eq:exposure}
\end{equation}

$f_{\text{rot}}$ penalizes missing or overdue rotation, $f_{\text{use}}$
penalizes zombie credentials that have been inactive for more than 180 days,
and $f_{\text{gov}}$ penalizes identities that lack governance metadata.

\subsection{AI-Amplification Factor ($A$)}

This is the central contribution of the paper. For non-agentic NHIs, $A = 1.0$
and the factor has no effect. For AI agents, we define:

\begin{equation}
  A(i) = \alpha(i) \times \beta(i) \times \gamma(i)
  \label{eq:amplification}
\end{equation}

\noindent\textbf{Autonomy weight} $\alpha$ reflects how much human oversight
exists:
\begin{equation}
  \alpha(i) = \begin{cases}
    1.0 & \text{human-in-the-loop} \\
    2.0 & \text{semi-autonomous} \\
    4.0 & \text{fully autonomous}
  \end{cases}
\end{equation}

\noindent\textbf{Tool blast radius} $\beta$ is the highest blast score among
all tools the agent can invoke (Table~\ref{tab:tool_blast}).

\noindent\textbf{Reversibility factor} $\gamma$ captures whether the agent
can take actions that cannot be undone:
\begin{equation}
  \gamma(i) = \begin{cases}
    2.5 & \text{if any irreversible tool is present} \\
    1.0 & \text{otherwise}
  \end{cases}
\end{equation}

The theoretical ceiling for $A$ is $4.0 \times 6.0 \times 2.5 = 60.0$,
reached by a fully autonomous agent with access to a financial transfer tool
and no human approval gate anywhere in its execution path.

\begin{table}[h]
\caption{Tool Blast Radius Scores (selected entries)}
\label{tab:tool_blast}
\centering
\begin{tabular}{lcc}
\toprule
\textbf{Tool} & \textbf{Blast Score} & \textbf{Irreversible} \\
\midrule
\texttt{transfer\_funds} & 6.0 & Yes \\
\texttt{bash / shell} & 5.5 & Yes \\
\texttt{execute\_code} & 5.0 & Yes \\
\texttt{deploy} & 5.0 & Yes \\
\texttt{delete\_record} & 4.5 & Yes \\
\texttt{delete\_file} & 4.0 & Yes \\
\texttt{update\_database} & 3.5 & No \\
\texttt{send\_email} & 3.0 & Yes \\
\texttt{call\_api} & 3.0 & No \\
\texttt{search\_web} & 1.0 & No \\
\bottomrule
\end{tabular}
\end{table}

\textbf{Why multiplicative?} Consider a fully autonomous agent ($\alpha = 4$)
with shell access ($\beta = 5.5$) that can run irreversible commands ($\gamma =
2.5$). Each tool invocation can chain into additional unauthorized actions at
machine speed, with no human checkpoint to break the cascade. Additive scoring
cannot represent this compounding dynamic; multiplication can.

\section{Implementation}

\subsection{Architecture}

AgentSentry is written in Python 3.10+ and organized around four modules:

\begin{itemize}
  \item \textbf{Scanners}: Cloud provider scanners covering AWS IAM, S3, and
        Lambda, plus a static AI agent analyzer built on Python's \texttt{ast}
        module.
  \item \textbf{Core}: Pydantic v2 data models, the scoring engine, and an NHI
        attack graph implemented with NetworkX.
  \item \textbf{Enrichment}: Integration with the CISA KEV catalog for
        real-time threat intelligence correlation.
  \item \textbf{Reporting}: A Click + Rich command-line interface, interactive
        HTML graph output via Pyvis, and JSON export for pipeline integration.
\end{itemize}

\subsection{AI Agent Static Analyzer}

The AI agent scanner uses Python AST parsing to extract agent definitions from
a codebase without running any code. This matters in practice: you can point
the tool at a production repository and get results without needing to install
its dependencies or set up its runtime environment.

The analyzer recognizes agent constructor patterns for LangChain, CrewAI, and
AutoGen. It extracts tool lists and autonomy configuration parameters, then
classifies each agent's autonomy level based on the presence of human approval
callbacks, \texttt{max\_iterations} settings, and tool reversibility flags.

\subsection{NHI Attack Graph}

The attack graph $G = (V, E)$ is a directed graph where:
\begin{itemize}
  \item $V = V_{\text{NHI}} \cup V_{\text{resource}}$: nodes for NHIs and
        the resources they can reach
  \item $E$: directed edges representing access relationships, weighted by
        traversal cost along the attack path
\end{itemize}

Blast radius for a compromised NHI $i$ is computed as:
\begin{equation}
  \text{BlastRadius}(i) = |\text{descendants}(i, G)| \times
  |\text{CrownJewels}(i, G)|
  \label{eq:blast}
\end{equation}

using NetworkX's \texttt{descendants()} function, which runs BFS over the
directed graph in $O(V + E)$ time.

\section{Evaluation}

\subsection{AWS Environment Validation}

We ran AgentSentry against a live AWS environment. The findings are summarized
in Table~\ref{tab:aws_results}.

\begin{table}[h]
\caption{AWS Scan Results}
\label{tab:aws_results}
\centering
\begin{tabular}{llccc}
\toprule
\textbf{Identity} & \textbf{Type} & \textbf{P} & \textbf{Risk} & \textbf{Level} \\
\midrule
test-ml-pipeline & IAM Role & 10.0 & 150.0 & CRITICAL \\
agentsentry-scanner & IAM Key & 10.0 & 150.0 & CRITICAL \\
\bottomrule
\end{tabular}
\end{table}

With CISA KEV enrichment turned on, both CRITICAL identities matched against
three ransomware-linked CVEs (CVE-2022-22954, CVE-2021-42287,
CVE-2021-42278). This connection matters: over-privileged IAM identities do
not exist in a vacuum. They operate inside an active exploitation landscape
where attackers know exactly which vulnerabilities to chain together.

\subsection{AI-Amplification Factor Validation}

Table~\ref{tab:agent_results} shows scoring behavior across three AI agent
archetypes drawn from the demo codebase.

\begin{table}[h]
\caption{AI Agent Scan Results}
\label{tab:agent_results}
\centering
\begin{tabular}{lccccc}
\toprule
\textbf{Agent} & \textbf{Autonomy} & $\alpha$ & $\beta$ & $\gamma$ & \textbf{Risk} \\
\midrule
CRM Agent & Full & 4.0 & 4.5 & 2.5 & CRITICAL (450) \\
Email Drafter & Semi & 2.0 & 3.0 & 2.5 & HIGH (99) \\
Research Agent & Semi & 2.0 & 1.5 & 1.0 & MEDIUM (33) \\
\bottomrule
\end{tabular}
\end{table}

Autonomy level turns out to be the dominant risk driver. The CRM agent holds no
cloud IAM permissions at all, yet it scores $4.5\times$ higher than the email
drafter purely because of its fully autonomous execution mode and access to
irreversible tools. This outcome validates the multiplicative structure of $A$
and, more importantly, it reveals something that privilege-only scoring
fundamentally cannot detect.

\subsection{Comparison with Existing Approaches}

Current CIEM tools score NHI risk on privilege alone. Under a pure privilege
model, the CRM Agent and a read-only IAM role with no attached policies both
score near zero. AgentSentry's inclusion of $A$ reveals a $450\times$
difference in actual risk between them. In environments with agentic AI
deployments, this is not an academic gap. It is the difference between finding
a critical vulnerability and missing it entirely.

\section{Discussion}

\subsection{Limitations}

Three limitations are worth being direct about. The static analyzer cannot
resolve tool lists that are constructed dynamically at runtime, for example,
tools loaded from a configuration file or database. The permission weight table
is hand-curated and may not generalize cleanly to custom IAM policies that
deviate from conventional naming patterns. Finally, the blast radius model does
not currently account for network segmentation controls that could limit lateral
movement even when permissions technically exist.

\subsection{Future Work}

Three directions seem most promising. Runtime interception of AI agent tool
invocations would allow more precise autonomy classification than static
analysis alone. Extending the scanner to Azure and GCP identity models would
improve coverage for organizations running multi-cloud environments. On the
commercial side, continuous monitoring, automated remediation workflows, and
audit-grade compliance reporting are planned as reserved features.

\section{Conclusion}

This paper introduced AgentSentry, an open-source framework for quantifying
risk in non-human identities and autonomous AI agents. The AI-Amplification
Factor ($A$) at the core of the framework formally captures something that
every practitioner working with agentic AI systems recognizes intuitively: an
autonomous agent running under a compromised identity is not just as dangerous
as that identity's permissions suggest. It is potentially far more dangerous,
because the agent can act at machine speed across every tool it controls,
without waiting for a human to notice. Our evaluation puts a number on that
intuition, showing $A$ increases composite risk scores by up to $60\times$
relative to non-agentic identities at identical privilege levels, and revealing
a class of critical vulnerabilities that existing privilege-only frameworks
will consistently miss.

AgentSentry is publicly available at
\url{https://github.com/Abhiram-ops/agent-sentry}.

\begin{thebibliography}{00}

\bibitem{crowdstrike2024}
CrowdStrike, ``2024 CrowdStrike Global Threat Report,'' CrowdStrike, Tech. Rep., 2024.

\bibitem{langchain2023}
H. Chase, ``LangChain,'' 2023. [Online]. Available: \url{https://github.com/langchain-ai/langchain}

\bibitem{crewai2024}
J. Moura, ``CrewAI: Framework for Orchestrating Role-Playing Autonomous AI Agents,'' 2024.
[Online]. Available: \url{https://github.com/joaomdmoura/crewAI}

\bibitem{autogen2023}
Q. Wu et al., ``AutoGen: Enabling Next-Gen LLM Applications via Multi-Agent Conversation,''
\textit{arXiv preprint arXiv:2308.08155}, 2023.

\bibitem{mitre2024}
MITRE Corporation, ``MITRE ATT\&CK Enterprise Matrix,'' 2024.
[Online]. Available: \url{https://attack.mitre.org}

\bibitem{owasp_llm2023}
OWASP, ``OWASP Top 10 for Large Language Model Applications,'' 2023.
[Online]. Available: \url{https://owasp.org/www-project-top-10-for-large-language-model-applications/}

\bibitem{perez2022}
F. Perez and I. Ribeiro, ``Ignore Previous Prompt: Attack Techniques for Language Models,''
\textit{arXiv preprint arXiv:2211.09527}, 2022.

\bibitem{greshake2023}
K. Greshake et al., ``Not What You've Signed Up For: Compromising Real-World LLM-Integrated
Applications with Indirect Prompt Injection,'' \textit{arXiv preprint arXiv:2302.12173}, 2023.

\bibitem{yi2023}
J. Yi et al., ``Benchmarking and Defending Against Indirect Prompt Injection Attacks on
Large Language Models,'' \textit{arXiv preprint arXiv:2312.14197}, 2023.

\bibitem{sheyner2002}
O. Sheyner et al., ``Automated Generation and Analysis of Attack Graphs,''
in \textit{Proc. IEEE Symposium on Security and Privacy}, 2002, pp. 273--284.

\bibitem{saltzer1975}
J. H. Saltzer and M. D. Schroeder, ``The Protection of Information in Computer Systems,''
\textit{Proc. IEEE}, vol. 63, no. 9, pp. 1278--1308, 1975.

\bibitem{aws_iam2024}
Amazon Web Services, ``AWS Identity and Access Management User Guide,'' 2024.
[Online]. Available: \url{https://docs.aws.amazon.com/IAM/latest/UserGuide/}

\bibitem{azure_identity2024}
Microsoft, ``Managed Identities for Azure Resources,'' 2024.
[Online]. Available: \url{https://learn.microsoft.com/en-us/azure/active-directory/managed-identities-azure-resources/}

\bibitem{k8s_sa2024}
Kubernetes, ``Configure Service Accounts for Pods,'' 2024.
[Online]. Available: \url{https://kubernetes.io/docs/tasks/configure-pod-container/configure-service-account/}

\bibitem{conjur2022}
CyberArk, ``Conjur Open Source,'' 2022.
[Online]. Available: \url{https://www.conjur.org}

\bibitem{workload_identity2023}
Google Cloud, ``Workload Identity Federation,'' 2023.
[Online]. Available: \url{https://cloud.google.com/iam/docs/workload-identity-federation}

\bibitem{access_analyzer}
Amazon Web Services, ``AWS IAM Access Analyzer,'' 2024.
[Online]. Available: \url{https://aws.amazon.com/iam/access-analyzer/}

\bibitem{orca2023}
Orca Security, ``Cloud Attack Path Analysis,'' 2023.
[Online]. Available: \url{https://orca.security}

\bibitem{wiz2023}
Wiz, ``Cloud Security Graph,'' 2023.
[Online]. Available: \url{https://www.wiz.io}

\end{thebibliography}

\end{document}
